%% =====================================================================
%%  B6-T-06-02 -- what B6 contributes to "The Two-Tellings Test"
%%  -------------------------------------------------------------------
%%  LAYER A: APPEND-ONLY. Written once at the entry's write, then left
%%  alone. If the source has more to say later, add a bullet -- do not
%%  rewrite. Material found ONLY here goes in the `unique` slot with
%%  @unique set to yes, which makes it undeletable (rule R10).
%% =====================================================================
%% @kind: source
%% @id: B6-T-06-02
%% @book: B6
%% @topic: T-06-02
%% @locator: ch. 9 (diversion, evasion, vagueness, rationalization, feigning innocence), pp. 112--124
%% @status: distilled
%% @unique: no
%% @integrated: yes
%% @updated: 2026-09-19

\dossier{B6}{T-06-02}{\bookshort{B6}}{ch. 9 (diversion, evasion, vagueness, rationalization, feigning innocence), pp. 112--124}

\begin{distilled}
  \item Manipulators are notorious for never giving a straight answer to a straight question: diversion changes the subject; evasion and vagueness produce answers that feel like answers and are not --- sometimes so subtle you think you have one when you have none.
  \item Rationalization is the grown-up version: plausible justifications delivered combined with other tactics (shaming, guilt) so the excuse and the pressure arrive together.
  \item Feigning innocence: the harm was unintentional, or never happened --- designed to make the target question their judgment and possibly their sanity; sometimes as little as a look of surprise or indignation.
  \item B6's detection rule for this family: you need a sensitive ear --- judge answers by whether they address the question asked, not by their fluency.
\end{distilled}
